\chapter{はじめに}

\section*{この本のねらい}

この本は、統計検定「データサイエンスエキスパート」の出題範囲を、
\textbf{高校数学と Python の文法だけ}を前提に学べるように書いた教科書です。
大学で学ぶ数学（行列やガンマ関数など）が必要になる所では、その場で最小限の説明を入れています。

章立ては、試験の出題範囲表に合わせています。
範囲表の「大項目」を部、「中項目」を章、「小項目」を節としました。
各節の最初には、範囲表に書かれているキーワードを示しています。

\section*{この本の読み方}

新しい用語が初めて出てくる所では、{\sffamily\gtfamily\bfseries 標本平均}（sample mean）のように太字で示し、英語の名前も添えます。
英語の名前は、英語の資料を読むときや、Python のライブラリで関数を探すときに役立ちます。
太字の用語は、巻末の索引にも載せています。

本文の中では、次のような囲みを使います。
定義・定理・例・例題には、章ごとの通し番号を付けます。
たとえば「定理 3.2」は、第 3 章の 2 番目の囲みという意味です。

\begingroup
\renewcommand{\thechapter}{3}  % 見本として「第3章」の番号を表示する

\begin{definition}
  言葉の意味を決めるものです。後の議論はすべてここから始まります。
\end{definition}

\begin{theorem}
  定義から導かれる、大事な性質です。この本では、定理を必ず\textbf{証明}（導出）します。
\end{theorem}

\begin{intuition}
  数式や定理が何を言っているのかを、言葉や図で説明します。
\end{intuition}

\begin{merit}
  その定理や数式を知っていると、何が楽になるのか、何ができるようになるのかを説明します。
  式を覚えるより先に、ここを読んでください。
\end{merit}

\begin{example}
  具体的な数値や場面で、定義や定理を確かめます。
\end{example}

\begin{problem}
  自分で手を動かして解いてみる問題です。
  \tcblower
  解答は、点線の下にあります。
\end{problem}
\endgroup

\begin{mathnote}
  高校数学を超える道具を、その場で必要なだけ説明します。
  くわしくは「数学基礎」の部で扱います。
\end{mathnote}

\begin{supplement}
  範囲表には書かれていませんが、後の内容を理解するために必要な話です。
\end{supplement}

\begin{reference}
  発展的な内容です。名前と大まかな意味を知っていれば十分です。
\end{reference}

\begin{exampoint}
  試験でよく問われる点や、まちがえやすい点です。
\end{exampoint}

Python のコードは次のように示します。
コードはすべて \texttt{numpy} などの基本的なライブラリだけで動くようにしています。

\begin{pycode}
import numpy as np
print(np.mean([1, 2, 3]))
\end{pycode}

\begin{pyoutput}
2.0
\end{pyoutput}

\section*{記号について}

\begin{center}
  \begin{tabular}{ll}
    \toprule
    記号 & 意味 \\
    \midrule
    $\Prob(A)$ & 事象 $A$ が起こる確率 \\
    $\E[X]$, $\V[X]$ & 確率変数 $X$ の期待値、分散 \\
    $\Cov(X,Y)$ & $X$ と $Y$ の共分散 \\
    $X \sim \Normal(\mu,\sigma^2)$ & $X$ が平均 $\mu$、分散 $\sigma^2$ の正規分布にしたがう \\
    $X_1,\dots,X_n \iid F$ & $X_1,\dots,X_n$ が互いに独立に、同じ分布 $F$ にしたがう \\
    $\bX$ & 標本平均 $\frac{1}{n}\sum_{i=1}^n X_i$ \\
    $\estim{\theta}$ & パラメータ $\theta$ の推定量 \\
    \bottomrule
  \end{tabular}
\end{center}
